
\documentclass{article}
\usepackage{colortbl}
\usepackage{makecell}
\usepackage{multirow}
\usepackage{supertabular}

\begin{document}

\newcounter{utterance}

\centering \large Interaction Transcript for game `codenames', experiment `unambiguous', episode 4 with reflexion{-}Qwen3.5{-}27B.
\vspace{24pt}

{ \footnotesize  \setcounter{utterance}{1}
\setlength{\tabcolsep}{0pt}
\begin{supertabular}{c@{$\;$}|p{.15\linewidth}@{}p{.15\linewidth}p{.15\linewidth}p{.15\linewidth}p{.15\linewidth}p{.15\linewidth}}
    \# & $\;$A & \multicolumn{4}{c}{Game Master} & $\;\:$B\\
    \hline

    \theutterance \stepcounter{utterance}  
    & & & \multicolumn{2}{p{0.3\linewidth}}{
        \cellcolor[rgb]{0.9,0.9,0.9}{
            \makecell[{{p{\linewidth}}}]{
                \texttt{\tiny{[GM$|$GM]}}
                {'hidden': {'team': ['drugstore'], 'innocent': ['stewardess', 'blimp', 'appetizer', 'pharmacist', 'prom', 'geisha', 'shaman'], 'opponent': [], 'assassin': ['cowardice']}, 'revealed': {'team': {'team': [], 'innocent': [], 'opponent': [], 'assassin': []}, 'opponent': {'team': [], 'innocent': [], 'opponent': [], 'assassin': []}}}
            }
        }
    }
    & & \\ \\

    \theutterance \stepcounter{utterance}  
    & & \multicolumn{4}{p{0.6\linewidth}}{
        \cellcolor[rgb]{0.9,0.9,0.9}{
            \makecell[{{p{\linewidth}}}]{
                \texttt{\tiny{[P1$\langle$GM]}}
                \texttt{You have attempted this same task before and did not succeed. These are the plans you wrote after those attempts.} \\
\texttt{Trial \#1: New plan: My previous strategy failed because I prioritized safety over efficiency, giving clues for only one or two words at a time. This allowed the opponent to catch up and win before I could clear my list. Specifically, I missed opportunities to group words with broader associations: "bathtub" and "drugstore" could have been linked via "store" or "bath" (if "bath" didn't risk the kill word, though "pharmacy" was the fatal error as it contained the root of "pharmacist" in the distractors, or simply was too specific). More critically, I should have grouped "serenity," "prairie," and "daybreak" under a broader concept like "peace" or "morning" earlier. In the future, I will aggressively group 3{-}4 words per turn using broader, high{-}level semantic categories (e.g., "nature" for prairie/snowman/daybreak, "building" for apartment/stairway/bathtub) to maximize the number of words guessed per turn, accepting a slightly higher risk of a wrong guess to ensure I finish my list before the opponent. I must also strictly avoid clues that are substrings of distractor words (like "pharmacy" for "pharmacist") or too close to the kill word.} \\
\\ 
\texttt{Trial \#2: New plan: My previous attempt failed because I chose "pharmacy" as a clue, which is semantically too close to the distractor "pharmacist," likely causing my partner to guess the wrong word and end our turn. Additionally, I only targeted one word, which is inefficient. Next time, I will select a broader, safer clue that avoids substrings or direct synonyms of distractors and kill words. For "drugstore," I will use a clue like "shop" or "store" to clearly indicate the location without triggering the "pharmacist" distractor. I will also ensure the clue relates to as many of my team's words as possible in a single turn to maximize efficiency.} \\
\\ 
\texttt{Trial \#3: New plan: My previous attempt failed because I was overly conservative, targeting only one word ("drugstore") with a clue ("SHOP") that, while safe, did not maximize the turn's potential. Although the current list only contains one word, the strategy of limiting targets to a single item is inefficient in a game where speed is critical. In future turns with multiple words, I will aggressively group 3{-}4 words under a single, broad semantic category (e.g., "place," "item," "action") to clear the board faster. I will strictly ensure the clue is not a substring of any distractor (like avoiding "pharm" for "pharmacist") or the kill word. For this specific turn, since only "drugstore" exists, the clue "SHOP" was technically correct for the single target, but the failure implies the game state required a different approach or the opponent won due to my slow pace in previous turns; thus, I will prioritize high{-}volume grouping in all future turns to prevent the opponent from catching up.} \\
\\ 
\\ 
\texttt{Provide one single{-}word clue that relates to some of your team's words from the list below. You can choose to how many and to which words your clue relates to.} \\
\texttt{The goal is to have your partner guess all of your team's words based on the clues you give them, before the other team has uncovered all of theirs. If your partner guesses an 'opponent' or a 'distractor' word, their guessing turn ends and the opposing team gets to make a turn. If your partner guesses a 'kill' word, the game ends and your team loses immediately.} \\
\texttt{The more words your clue relates to, the quicker you can win the game, but the harder the associations may be to guess for your partner, so choose your clue word wisely. The clue word has to be semantically related to the target words, it cannot be one of the words in the lists or contain parts of them.} \\
\\ 
\texttt{Always give your single{-}word clue and your comma{-}separated list of related target words in the following format and make your answers as short as possible, never include any other text than is required in this form:} \\
\\ 
\texttt{CLUE: <WORD>} \\
\texttt{TARGETS: <TARGETS>} \\
\\ 
\texttt{Your team words are: drugstore.} \\
\texttt{Your opponent's team words are: .} \\
\texttt{Distractor words are: stewardess, blimp, appetizer, pharmacist, prom, geisha, shaman.} \\
\texttt{The kill words are: cowardice.} \\
            }
        }
    }
    & \\ \\

    \theutterance \stepcounter{utterance}  
    & \multicolumn{4}{p{0.6\linewidth}}{
        \cellcolor[rgb]{0.9,0.9,0.9}{
            \makecell[{{p{\linewidth}}}]{
                \texttt{\tiny{[P1$\rangle$GM]}}
                \texttt{CLUE: SHOP} \\
\texttt{TARGETS: drugstore} \\
            }
        }
    }
    & & \\ \\

    \theutterance \stepcounter{utterance}  
    & & & \multicolumn{2}{p{0.3\linewidth}}{
        \cellcolor[rgb]{0.9,0.9,0.9}{
            \makecell[{{p{\linewidth}}}]{
                \texttt{\tiny{[GM$|$GM]}}
                \texttt{shop} \\
            }
        }
    }
    & & \\ \\

    \theutterance \stepcounter{utterance}  
    & & & \multicolumn{2}{p{0.3\linewidth}}{
        \cellcolor[rgb]{0.9,0.9,0.9}{
            \makecell[{{p{\linewidth}}}]{
                \texttt{\tiny{[GM$|$GM]}}
                ['drugstore']
            }
        }
    }
    & & \\ \\

    \theutterance \stepcounter{utterance}  
    & & & \multicolumn{2}{p{0.3\linewidth}}{
        \cellcolor[rgb]{0.9,0.9,0.9}{
            \makecell[{{p{\linewidth}}}]{
                \texttt{\tiny{[GM$|$GM]}}
                {'word': 'drugstore', 'assignment': 'team'}
            }
        }
    }
    & & \\ \\

    \theutterance \stepcounter{utterance}  
    & & & \multicolumn{2}{p{0.3\linewidth}}{
        \cellcolor[rgb]{0.9,0.9,0.9}{
            \makecell[{{p{\linewidth}}}]{
                \texttt{\tiny{[GM$|$GM]}}
                \texttt{opponent has won} \\
            }
        }
    }
    & & \\ \\

    \theutterance \stepcounter{utterance}  
    & \multicolumn{4}{p{0.6\linewidth}}{
        \cellcolor[rgb]{0.95,0.95,0.95}{
            \makecell[{{p{\linewidth}}}]{
                \texttt{\tiny{[P2$\rightarrow$P1]}}
                \texttt{Post{-}episode Reflexion (not sent to GM):} \\
\texttt{New plan: My previous attempt failed because I provided a clue ("SHOP") that, while semantically related to "drugstore," was too generic and likely triggered the distractor "pharmacist" or failed to distinguish the target clearly enough in a competitive context. The clue "SHOP" is a broad category that could easily lead a partner to guess "pharmacist" (a person who works in a shop) or other distractors like "stewardess" (if associated with a shop on a plane) or "geisha" (if associated with a shop selling traditional items), causing the turn to end prematurely. To fix this, I must choose a clue that is uniquely and strongly associated with "drugstore" while explicitly avoiding any semantic overlap with the distractors (stewardess, blimp, appetizer, pharmacist, prom, geisha, shaman) and the kill word (cowardice). Instead of "SHOP," I will use a more specific functional clue like "MEDS" or "RX" (if allowed as a single word) or "COUNTER" which strongly implies a drugstore but is less likely to be confused with a person (pharmacist) or other unrelated items. If "MEDS" is too risky, "PRESCRIPTION" is too long, so "COUNTER" or "BOTTLE" might be safer. However, given the distractors, "PHARMACY" is forbidden (substring of pharmacist), so I will use "MEDICINE" or "DRUGS" (if not a substring of the target, but "drugstore" contains "drug", so "DRUGS" is risky). The best specific clue that avoids "pharmacist" and other distractors is likely "COUNTER" or "BOTTLE". I will select "COUNTER" as it is a specific part of a drugstore not shared by the other distractors, ensuring the partner guesses "drugstore" and not "pharmacist". I will also ensure that in future turns with multiple words, I group them under a single, precise clue that avoids all distractors and kill words, rather than playing it safe with one word at a time.} \\
            }
        }
    }
    & & \\ \\

\end{supertabular}
}

\end{document}
